\section{Experiments}\label{sec:experiments}

All experiments are small by design: CPU training, one geometry family, a small trunk. They isolate one design choice at a time
and do not probe the accuracy ceiling of the method. Frequency errors of one to two percent therefore reflect an
under-trained model, and only comparisons between arms that share model, family and evaluation protocol are evidence.

\paragraph{Setup.}
% src: docs/34 §3, §4, §7, §8.2; configs/field_3d.yaml; src/training/lightning_module.py (custom_cosine, AdamW)
Everything ran on a shared four-core cloud VM without a GPU, often next to other jobs, so wall-clock times scatter noticeably.
The equal-cost experiment used one compute thread per run, the earlier studies two. The model is that of \cref{sec:method} with embedding width~64,
two attention blocks, four heads, basis width $\nbasis=12$, 32 random Fourier features and $\nmodes=4$ modes. It is trained with
AdamW~\cite{loshchilov2019decoupled} in PyTorch~\cite{paszke2019pytorch}: learning rate $10^{-3}$, weight decay $10^{-4}$,
gradient-norm clipping at~1, batch size~1, no dropout, no augmentation. The learning rate is constant for the first ten epochs and then
decays with a cosine to $10^{-6}$ at the last epoch (it never decays in a run of ten epochs or fewer). The loss is the span loss in the
mass-plus-energy norm plus the self-supervised, orthogonality and frequency terms of \cref{sec:method} at default weights; the
figure-of-merit loss is off for the field model.

\emph{Data.} All geometries are of the family \texttt{pillbox\_pipes} (pillbox cavities with pipes), minimum-fillet parameter $0.12$,
drawn from a fixed-seed Sobol sequence, so identifier~$i$ is the same geometry in the LF and HF sets (checked on four samples).
Training sets: identifiers 0--7 (HF, 8 geometries), 0--23 (LF, 24), and 0--23 together with 30--45 (LF, 40). The six test geometries are
identifiers 24--29 with HF labels; their model meshes have 4.8k--10.0k tetrahedra, and the LF label meshes coincide with the HF model meshes (\cref{sec:multifidelity}).

\emph{Metrics.} The frequency error of mode~$k$ is $f_k^{\mathrm{model}}/f_k^{\HF}-1$ against the stored HF label; where the evaluated space differs
from the label space, and to test \cref{prop:upper-bound}, it is also formed against the finite-element eigenvalue of the evaluated space, $f_k^{\mathrm{model}}/f_k^{\mathrm{FE}}-1$.
The field error is the relative $L^2$ error in the mass norm against the finite-element mode of the evaluated space, minimised over the sign for an isolated mode,
and the relative distance between target and predicted subspace inside a cluster of near-degenerate modes. It is averaged over modes 0--2: the fourth mode is not yet learned
(field error $0.07$--$1.3$ in the first study below; cause not investigated). Models are evaluated from the last checkpoint, without checkpoint selection, on the stored mesh in the order-2 space.

\subsection{Representation accuracy}\label{sec:07-floor}
% src: docs/28 §3
\Cref{tab:07-overfit} overfits one curved pillbox (179 tetrahedra, width-64 two-block trunk, $\nbasis=12$, $\nmodes=4$). With vertex tokens and barycentric interpolation
the eigenvalue error saturates near $10^{-4}$ (worse at order~3); with the quadrature-point tokens of \cref{sec:method} it reaches $10^{-6}$--$9\times10^{-6}$ in fewer steps.
The lowest-order Whitney (N0) model is limited to a relative frequency error of $2.4\times10^{-3}$ on similar cavities (a different sample). A single-sample fit measures representational capacity, not generalisation.

\begin{table}[t]
\centering\small
\caption{Single-sample overfit on a curved pillbox: relative eigenvalue error $\theta_k/\lambda^{\mathrm{label}}_k-1$.}
\label{tab:07-overfit}
% src: docs/28 §3
\begin{tabular}{lccrc}
\toprule
Tokens & order $p$ & DOF & steps & eigenvalue error\\
\midrule
vertices + barycentric interpolation & 2 & 2280 & 2000 & $0.8$--$1.1\times10^{-4}$\\
vertices + barycentric interpolation & 3 & 5480 & 2000 & $2.3$--$3.0\times10^{-4}$\\
quadrature points & 2 & 2280 & 1250 & $\mathbf{1\text{--}9\times10^{-6}}$\\
\bottomrule
\end{tabular}
\end{table}

\subsection{Equal-cost comparison of low- and high-fidelity labels}\label{sec:07-equal-cost}
% src: docs/34 §8.1, §8.2
This is the main experiment. All arms share model, family, test set, optimiser and 600 optimisation steps; they differ in the labels.
Arm~A trains on the 8 HF geometries (each listed five times per epoch; 15 epochs of 40 items). Arm~B2 uses 24 LF geometries (25 epochs of 24), about $20\%$ cheaper in label time than A.
Arm~C2 uses 40 LF geometries (15 epochs of 40, the schedule of A) at $1.4\times$ the label cost of A. Label cost is the measured time for mesh generation plus labelling on the shared VM
(two workers, one thread each, next to a training job); the HF/LF cost ratio of the same eight meshes is $4.3$. Arm~D continues C2 for 200 steps on the 8 HF geometries at learning rate $3\times10^{-4}$.

\begin{table}[t]
\centering\small
\caption{Equal-cost comparison on the six HF test geometries (identifiers 24--29). $|\Delta f|$: absolute relative frequency error against the HF label, summarised over geometries and modes; field: mass-norm relative $L^2$ error averaged over modes 0--2. One training run per arm. Arm~D has no valid control (see text).}
\label{tab:equal-cost}
% src: docs/34 §8.2 (table)
\setlength{\tabcolsep}{5pt}
\begin{tabular}{llrrrrrc}
\toprule
Arm & Training labels & Label cost [s] & Steps & \multicolumn{3}{c}{$|\Delta f|$ [\%]} & field $\relLtwo$\\
\cmidrule(lr){5-7}
 & & & & mean & median & max & \\
\midrule
A  & HF, 8 geometries  & 1242 & 600 & 1.66 & 1.11 & 8.11 & 0.050\\
B2 & LF, 24 geometries & 997  & 600 & 1.55 & 1.06 & 6.57 & 0.046\\
C2 & LF, 40 geometries & 1704 & 600 & \textbf{1.37} & \textbf{1.00} & \textbf{5.30} & \textbf{0.045}\\
D  & C2, then HF 8 fine-tune & 1704 + 1242 & 600 + 200 & 1.13 & 0.82 & 4.41 & 0.039\\
\bottomrule
\end{tabular}
\end{table}

\paragraph{Result.} At equal or lower label cost, training on LF labels was at least as accurate on the held-out HF geometries as training on HF labels, and more LF geometries helped slightly (\cref{tab:equal-cost}).
B2 costs $20\%$ less than A and has a lower mean error ($1.55\%$ against $1.66\%$), but this difference comes mostly from the outlier geometry~27; without it the means are $1.11\%$ (A), $1.14\%$ (B2) and $1.06\%$ (C2).
% src: docs/34 §8.2 (means without geometry 27: HF 8 1.11 %, LF 24 1.14 %, LF 40 1.06 %; C2 better than A on 3 of 6 geometries)
C2, at $1.4\times$ the cost, is lower than A on three of the six geometries (24, 27, 29) and $0.01$--$0.06$ percentage points behind on the other three (per-geometry means for identifiers 24--29, A against C2, in \%: $0.74/0.66$, $1.18/1.23$, $1.24/1.25$, $4.44/2.93$, $0.92/0.98$, $1.46/1.17$).
Each arm is a single training run and the differences are small against the spread over geometries, so we do not rank the arms; the supported statement is that LF labels did not hurt. Geometry~27 is the worst in every arm ($2.9$--$4.4\%$); we did not investigate why.

\paragraph{LF-to-HF gap.} The predictions of C2 have a mean error of $1.35\%$ against the LF labels of the test meshes and $1.37\%$ against the HF labels: the model trained on LF data is $0.02$ percentage points
worse on HF targets than on LF targets. This matches \cref{sec:multifidelity}: on identical meshes the LF and HF frequencies differ by $0.004$--$0.047\%$ and the HF mode lies outside the six-mode LF space by $0.08$--$0.32\%$ in the mass norm (modes 0--3);
the model error is 10 to 50 times larger.
% src: docs/34 §8.1, §8.2

\paragraph{Ritz bound.} In every evaluation of this experiment the predicted eigenvalues lie above the finite-element eigenvalue of the evaluated space, with smallest margin $+0.47\%$ (\cref{prop:upper-bound}).
% src: docs/34 §8.2; §4.1 (+0.006 %), §4.2 (+0.009 %) for the earlier studies

\paragraph{Fine-tuning and mixed training are inconclusive.} Arm~D improves on C2 (mean error $1.37\%\to1.13\%$), but the gain cannot be attributed to the HF labels. The control (same start, same 200 steps, LF data only) was lost when the container restarted before its test evaluation;
on the three validation geometries it had a mean absolute frequency error of $0.031$\,GHz and a field error of $0.046$, against $0.029$\,GHz and $0.052$ for D, a mixed signal. Two further runs are excluded: a mixed LF-40 plus HF-8 run with $8\times80$ steps never entered its decay phase
(eight epochs is fewer than the ten-epoch constant phase) and reached $1.95\%$, and a mixed LF-32 plus HF-8 run with the same schedule was lost on restart at epoch~6.
% src: docs/34 §8.2 (invalid and lost arms)
We therefore make no claim about fine-tuning or mixed training; they need a corrected schedule, a control arm and larger scale.
An earlier, smaller study (14 LF training geometries, evaluation on four HF geometries; \cref{sec:A07-first}) shows the same picture: the frequency errors on HF geometries were $0.63$--$3.02\%$ and all positive, the discrete order-2 eigenvalue on the model mesh lies within $0.033\%$ of the HF labels (the floor of the model space), so the error is the model's, and a fine-tuning test with a proper control gave no gain from two HF geometries (the improvement $1.8\%\to1.5\%$ came from the extra epochs).
% src: docs/34 §4.1, §4.3

\subsection{Transfer to other spaces}\label{sec:07-transfer}
% src: docs/34 §4.2
The Ritz layer uses only $\Kmat,\Mmat,\Gmat,\Kp$ of the space it is given, so the trained weights can be evaluated in another conforming space without retraining. \Cref{tab:07-transfer} does this for the model of the first study:
the native order-2 space, order~3 on the same mesh, and, for geometry~5, order~2 on the once uniformly refined mesh.

\begin{table}[t]
\centering\small
\caption{Same weights in different spaces. $\Delta f$: model against the HF label, signed percent, modes 0--3; field $\relLtwo$: modes 0--2 against the finite-element modes of the space. Times in seconds (two threads): operator assembly plus model forward, against a finite-element solve of the same space.}
\label{tab:07-transfer}
% src: docs/34 §4.2
\setlength{\tabcolsep}{3.5pt}
\begin{tabular}{llrlcrr}
\toprule
Id & Space & DOF & $\Delta f$ [\%] & field $\relLtwo$ & ops + fwd & FE solve\\
\midrule
3 & $p{=}2$ native   & 74\,301  & $+1.61/+1.05/+1.16/+2.51$ & .037/.064/.072 & 3.8 + 1.3  & 11.5\\
3 & $p{=}3$          & 186\,576 & $+2.29/+1.39/+1.51/+3.23$ & .038/.065/.073 & 12.6 + 5.1 & 48.1\\
5 & $p{=}2$ native   & 64\,965  & $+1.29/+0.63/+0.68/+1.89$ & .026/.043/.037 & 1.7 + 1.2  & 6.4\\
5 & $p{=}3$          & 163\,176 & $+1.88/+0.88/+0.94/+2.53$ & .027/.044/.037 & 8.4 + 4.4  & 27.6\\
5 & $p{=}2$, refined ($8\times$ tets) & 489\,747 & $+1.89/+0.86/+0.89/+2.29$ & .027/.043/.037 & 32.3 + 22.0 & 192\\
\bottomrule
\end{tabular}
\end{table}

The bound holds in all spaces (smallest margin $+0.009\%$ to the eigenvalue of the evaluated space), but accuracy does not improve: the frequency error is $0.2$--$0.6$ percentage points \emph{worse} than natively and the field error is unchanged.
A richer space enlarges the admissible fields, but the learned basis $\psi_j$ is the same smooth field, so the error comes from the basis and not from the space; the token density (24 quadrature points per tetrahedron at order~3 against 14 at order~2) also differs from training, a small distribution shift.
Transfer thus establishes validity (guarantee and absence of spurious modes survive a change of space), not accuracy.

\subsection{Inference cost}\label{sec:07-cost}
% src: docs/34 §4.2, §6, §8.2
At these sizes a direct finite-element solve is of the same order of magnitude as the model. In the native order-2 space, operators plus forward pass took $5.1$ and $2.9$\,s on geometries 3 and 5 against $11.5$ and $6.4$\,s for the solve (\cref{tab:07-transfer});
in the equal-cost study (one thread) they took $5$\,s and $3.3$\,s against $18$\,s. The model is therefore about a factor of two faster (ratios computed from these timings), $2.2$--$2.7$ at order~3 and $3.5$ on the refined mesh ($54$\,s against $192$\,s),
while its error is $1$--$2\%$ and the solve is exact up to the discretisation error of its space. We claim no speed advantage at this scale. Any benefit lies in larger spaces, on a GPU where the network share would shrink (not measured), or in warm-starting an iterative eigensolver from the model subspace (not implemented).

\subsection{Training-path engineering}\label{sec:07-engineering}
% src: docs/31 §3.6, §3.7; docs/28 §4
Training on curved high-order meshes is dominated by work that does not depend on the learned parameters. Removing repeated work (stored wall features, a transpose prepared once per geometry for the sparse products, preallocated buffers in the conjugate-gradient projection) cut the single-thread
training step of the 8931-tetrahedron pillbox from $40.7$ to $32.4$\,s ($-20\%$) and the dataset loading time by $37$--$41\%$; activation checkpointing in the trunk reduced the peak memory of a 13k-tetrahedron step from $9.9$ to $3.8$\,GB with unchanged step time on CPU.
The loss, the input tensors and all 76 parameter gradients are bit-identical to the previous code on two geometries, and the gradients with checkpointing are identical as well. The measurements, including a smaller gain on a smaller cell and the changes that were measured and not adopted, are in \cref{sec:A07-engineering} (\cref{tab:A07-engineering}); no GPU timing was made.
